%-------------------------------------------------------------------------------
% Chapter 2: Background and Related Work
%-------------------------------------------------------------------------------

\chapter{Background and Related Work}\label{chap:research}

To properly contextualize the operational environment of this research, it is necessary to examine the foundational concepts of continuous black-box optimization alongside the algorithmic vulnerabilities introduced by stochastic noise. Building upon these theoretical foundations, recent advancements in Large Language Models (LLMs) have driven a paradigm shift toward Automated Algorithm Design (AAD), enabling autonomous agents to synthesize executable search strategies. This chapter establishes the mathematical background of continuous optimization under uncertainty and surveys the state-of-the-art literature in LLM-driven code evolution.

%-------------------------------------------------------------------------------
\section{Theoretical Background}\label{sec:bg_theory}

Continuous black-box optimization focuses on identifying the global minimum of an unknown, real-valued objective function:
\begin{equation}\label{eq:bbob_min}
    \min_{x \in \mathcal{S} \subseteq \mathbb{R}^D} f(x)
\end{equation}
where $\mathcal{S}$ represents a bounded $D$-dimensional search space (typically $[-5, 5]^D$), and the analytical gradient $\nabla f(x)$—alongside any internal structural information regarding the function—is completely unavailable, leaving point sampling as the sole mechanism for information retrieval. In standard literature, algorithms designed for these problems are systematically evaluated against the deterministic Black-Box Optimization Benchmark (BBOB) suite \cite{Hansen2010}. Standardized within the IOHexperimenter and COCO platforms, the BBOB testbed comprises 24 synthetic mathematical functions designed to emulate distinct topological traps encountered in real-world engineering. These functions are categorized into five structural problem groups, as detailed in Table~\ref{tab:bbob_groups}.

\begin{table}[htbp]
\centering
\caption{Taxonomy of the 24 Continuous Black-Box Optimization Benchmark (BBOB) Functions \cite{Hansen2010}.}
\label{tab:bbob_groups}
\small
\begin{tabularx}{\textwidth}{c c l X}
\toprule
\textbf{Group} & \textbf{Function Indices} & \textbf{Landscape Property} & \textbf{Core Topological Challenge} \\ 
\midrule
1 & $f_1 - f_5$   & Separable Functions               & Dimensions are orthogonal; can be optimized independently along axes. \\
2 & $f_6 - f_9$   & Low/Moderate Conditioning        & Unimodal landscapes with smooth, single-valley trajectories. \\
3 & $f_{10} - f_{14}$ & High Conditioning            & Steep, narrow ill-conditioned valleys requiring precise coordinate rotations. \\
4 & $f_{15} - f_{19}$ & Multi-modal w/ Global Structure & Multiple deceptive local minima superimposed on an overarching global slope. \\
5 & $f_{20} - f_{24}$ & Multi-modal w/ Weak Structure   & Highly deceptive, randomized local optima lacking strong global gradients. \\ 
\bottomrule
\end{tabularx}
\end{table}

Crucially, the BBOB framework standardized the use of empirical cumulative distribution functions (ECDFs) to characterize algorithm performance, measuring runtime by the number of function evaluations required to achieve a specific target precision ($\Delta f_t = f(x) - f_{\text{opt}} \le 10^{-8}$). Over recent decades, human engineering has yielded robust heuristics to navigate these deterministic spaces, most notably the Nelder-Mead simplex algorithm \cite{Nelder1965}, Differential Evolution (DE) \cite{Storn1997}, Particle Swarm Optimization (PSO) \cite{Kennedy1995}, and the Covariance Matrix Adaptation Evolution Strategy (CMA-ES) \cite{Hansen2001}. Empirical benchmarks on the noiseless BBOB testbed consistently demonstrate that CMA-ES represents the state-of-the-art baseline for continuous optimization, outperforming local search methods like NEWUOA \cite{Powell2006} by multiple orders of magnitude on highly conditioned, non-separable topologies due to its robust covariance matrix update mechanics \cite{Hansen2010}.

While human-engineered algorithms excel on deterministic functions, practical industrial applications—such as electro-chemical battery cell parameter extraction \cite{Plett2004, Xiong2018} and structural finite-element model updating \cite{Mottershead1993, Marwala2010}—are inherently non-deterministic. Objective evaluations in these physical environments are derived directly from hardware instrumentation, telemetry apparatuses, or complex numerical simulations, rendering them persistently corrupted by thermal drift, electromagnetic interference, sensor quantization limits, and environmental turbulence \cite{Plett2004, Xiong2018}. As established in the foundational taxonomic framework on uncertain optimization environments by Jin and Branke \cite{Jin2005}, stochastic bottlenecks in computational search can be classified by whether uncertainty disrupts the parameter space through variable deviations post-optimization, the time space via dynamic topological drift, the surrogate space through approximation errors, or the fitness space. This thesis isolates noise acting directly within the fitness space, where the true, time-invariant system response $f_{\text{clean}}(x)$ is masked by an additive stochastic measurement error component $z$:
\begin{equation}\label{eq:noisy_evaluation}
    f_{\text{noisy}}(x) = f_{\text{clean}}(x) + z
\end{equation}
rendering the observed objective value a random variable centered around the underlying true fitness, which directly distorts candidate evaluations during the optimization process \cite{Jin2005}.

Standard continuous optimization testbeds model these environmental disruptions through three primary statistical noise distributions \cite{Hansen2009Noisy}:
\begin{enumerate}
    \item \textbf{Additive Uniform Noise:} Models digital sensor quantization errors and signal truncation limits:
    \begin{equation}
        z \sim \mathcal{U}(-a, a), \quad p(z) = \frac{1}{2a} \quad \text{for } z \in [-a, a]
    \end{equation}
    
    \item \textbf{Additive Cauchy Noise:} A heavy-tailed distribution modeling severe, intermittent physical sensor dropouts or hardware spikes:
    \begin{equation}
        p(z) = \frac{1}{\pi \gamma \left[1 + \left(\frac{z}{\gamma}\right)^2\right]}
    \end{equation}
    
    \item \textbf{Additive Gaussian Noise:} Models background thermal, electronic, and environmental white noise operating independently of signal magnitude \cite{Hansen2009Noisy}:
    \begin{equation}\label{eq:gaussian_noise}
        z \sim \mathcal{N}(0, \sigma^2), \quad p(z) = \frac{1}{\sigma \sqrt{2\pi}} \exp\left(-\frac{z^2}{2\sigma^2}\right)
    \end{equation}
\end{enumerate}

While the official BBOB-noisy suite \cite{Hansen2009Noisy} formalizes scale-invariant multiplicative log-normal distributions ($f \cdot \exp(\beta \mathcal{N}(0,1))$), physical hardware sensors exhibit constant-variance additive noise. Applying unscaled static additive noise across heterogeneous landscapes causes severe evaluation imbalances: a fixed noise level completely overwhelms flat separable functions while remaining mathematically imperceptible on steep, highly conditioned functions. Resolving this discrepancy requires engineering a normalized landscape-scaling mechanism, as detailed in Chapter 3.

In these turbulent configurations, conventional derivative-free heuristics collapse because selection operators in algorithms like CMA-ES rely on sorting population candidates by relative rank, which localized fitness noise systematically corrupts \cite{Jin2005, Miller1996}. Jin and Branke's survey confirmed that traditional human-engineered workarounds—such as static sample averaging or increasing population sizes—severely degrade optimization speed and over-consume function evaluation budgets due to the $\sqrt{N}$ sample size penalty \cite{Jin2005}. Consequently, a sub-optimal candidate may secure a dominant position purely due to a favorable statistical anomaly, deceiving the algorithm into adjusting its covariance matrices or search steps toward false gradients \cite{Jin2005, Hansen2009Noisy}. This vulnerability is magnified in classical gradient-based algorithms where minor statistical noise completely corrupts finite-difference derivative vectors, causing immediate algorithmic failure \cite{Sarafian2024}. Fixed hyperparameter configurations cannot dynamically react to local noise fluctuations, creating an active mandate to investigate whether autonomous logic generators can discover deep structural alternatives rather than relying on static computational averages.


%-------------------------------------------------------------------------------
\section{Related Work}\label{sec:related_work}


To eliminate this human bottleneck, the field is currently experiencing a profound paradigm shift towards Automated Algorithm Design (AAD). While early AAD frameworks relied on Genetic Programming \cite{Burke2013} or modular human-authored spaces \cite{vanRijn2016}, the immense size of the search spaces or strict boundaries frequently constrained discovery. This barrier has been recently overcome by the integration of Large Language Models, which were validated for strict mathematical discovery by Romera-Paredes et al. \cite{Romera2023}. Their framework, FunSearch, addresses the inherent tendency of LLMs toward hallucination by pairing the model with a systematic, closed-loop evaluator. By constraining the LLM to evolve the critical logical components of a program skeleton—rather than attempting to generate raw numerical solutions—the authors demonstrated that the model could synthesize novel, verifiable algorithmic constructions. Empirically, FunSearch achieved a historic milestone by discovering entirely new, large-sized caps in the cap-set problem that surpassed all existing lower bounds established by human mathematicians, alongside generating highly compact, reusable bin-packing heuristics that significantly outperformed traditional greedy algorithms across large-scale datasets \cite{Romera2023}. This landmark verification proves that foundation models possess the structural reasoning needed to invent mathematical scripts, paving the way to apply this exact closed-loop paradigm to real-valued exploration spaces.

Following this breakthrough, subsequent literature explored alternative modalities for utilizing LLMs in optimization pipelines. For instance, the EvoLLM framework \cite{Lange2024} established that an LLM can act directly as an evolutionary numerical sampler to predict coordinate distributions; however, this architecture lacks the interpretability of reusable code and incurs heavy prompt-token footprints. When benchmarked, EvoLLM demonstrated reasonable performance on a subset of eight continuous BBOB functions, but its validation was strictly restricted to a single 5-dimensional instance of each problem, highlighting severe scalability and generalizability limits \cite{Lange2024}. Conversely, frameworks such as the Evolution of Heuristics (EoH) \cite{Liu2024}, Algorithm Evolution using LLMs (AEL) \cite{AEL2024}, and Reflective Evolution (ReEvo) \cite{ReEvo2024} isolate the LLM's generative role to writing reusable code blocks. These approaches achieved superior performance and lower optimality gaps compared to traditional heuristics and deep reinforcement learning models, though their evaluations remained fundamentally confined to small, discrete instances of combinatorial tasks such as the Traveling Salesperson Problem (TSP) \cite{AEL2024, Liu2024}. 

Parallel advancements in language model multi-agent structures, most notably the Reflexion framework \cite{Shinn2023}, automated self-debugging architectures \cite{Chen2023}, and structural Abstract Syntax Tree (AST) feedback extensions like LLaMEA-SAGE \cite{LLaMEASAGE2026}, have confirmed that the introduction of text-based execution traces and code tracebacks back to the generator model drastically minimizes search space overhead. Specifically, Chen et al.\ proved that instructing a model to perform zero-human feedback "rubber duck debugging"—by generating line-by-line natural language code explanations alongside unit test execution checks—boosts functional accuracy by up to 12\% on programmatic translation tasks and improves output success on highly complex logical problems by 9\% while matching the performance of baselines demanding a 10$\times$ higher sampling budget \cite{Chen2023}. This extreme sample efficiency provides the structural justification for utilizing automated, text-based runtime execution feedback as a primary driver for heuristic optimization.

Building directly upon this proven capacity for executable code generation and automated self-correction, this thesis extends the Large Language Model Evolutionary Algorithm (LLaMEA) framework \cite{LLaMEA2024}. Engineered specifically to automate the discovery of continuous black-box optimization heuristics, LLaMEA operationalizes this closed-loop concept by structuring an LLM within a $1+1$ Evolution Strategy. The framework iteratively prompts the model to generate complete, reusable Python classes, executes that code against target mathematical problems via the IOHexperimenter benchmarking suite, and subsequently feeds both runtime error tracebacks and aggregated performance metrics—specifically the Area Over the Convergence Curve (AOCC)—back into the LLM's context window to drive self-debugging and refinement. Mathematically, the AOCC metric evaluated over a budget of $B$ evaluations on target precision $\Delta f_t$ is defined as:
\begin{equation}
    \text{AOCC} = \frac{1}{B} \sum_{k=1}^{B} \left( 1 - \frac{\log_{10}(\max(\Delta f(x_k), 10^{-8}))}{\log_{10}(\Delta f_{\text{init}})} \right)
\end{equation}

Empirically, foundational experiments proved that LLaMEA-generated algorithms successfully discovered novel heuristic logic that outcompeted state-of-the-art metaheuristics, including vanilla CMA-ES and Differential Evolution, on the 5-dimensional continuous BBOB suite, while maintaining highly competitive generalizability when zero-shot scaled to unseen 10-dimensional and 20-dimensional instances \cite{LLaMEA2024}. However, existing applications of the framework are strictly bounded to clean, noiseless functions, representing a critical, unexamined boundary ceiling. If an automated code class optimized strictly under deterministic expectations experiences persistent statistical noise in the fitness space, its selection tracking will inevitably decay due to ranking corruption \cite{Jin2005, Hansen2009Noisy}. By systematically injecting a controlled statistical noise profile into the evaluation sandbox, this thesis evaluates an extension of the LLaMEA architecture—investigating whether closed-loop evolutionary feedback can drive a language model to discover specialized, noise-resilient optimization heuristics.


% ==============================================================================
% EXPANSION ROADMAP FOR THE 10+ PAGE DEEP DIVE (TO BE CODED AS CONTENT LATER)
% ==============================================================================
%
% % FUTURE SECTION 1: Deep Dive on Traditional Paradigms (To expand Para 1)
% % - Placeholder for formal mathematical definitions of the continuous problem f(x).
% % - Placeholder for a LaTeX table mapping out the 24 BBOB functions grouped into their 5 landscape classes.
% % - Placeholder for a sub-section on ECDF aggregation physics and target precision math.
%
% % FUTURE SECTION 2: Mathematical Analysis of Stochastic Interference (To expand Para 2)
% % - Placeholder to insert LaTeX math blocks with full PDFs (Probability Density Functions) for Uniform, Gaussian, and Cauchy distributions.
% % - Placeholder to detail the math of Hessian corruption and gradient-chasing failures in CMA-ES.
% % - Placeholder for an evaluation of standard human mitigations (e.g., re-sampling coefficients, margin thresholds).
%
% % FUTURE SECTION 3: The Mechanical Evolution of Early AAD Frameworks (To expand Para 3)
% % - Placeholder to contrast Genetic Programming syntax constraints (AST bloat) with LLM programmatic generations.
% % - Placeholder to analyze the FunSearch database structure and "Program Skeleton" parsing constraints.
%
% % FUTURE SECTION 4: Comparative Breakdown of Alternative Modalities (To expand Para 4)
% % - Placeholder for a comprehensive comparison matrix table matching: EvoLLM vs. EoH vs. AEL vs. LLaMEA.
% % - Placeholder to map out token economics and runtime black-box costs of numerical sampling methods.
% % - Placeholder to break down Reflexion multi-agent tracking schemas and AST validation loops.
%
% % FUTURE SECTION 5: Operational Limits and Boundaries of LLaMEA 2024 (To expand Para 5)
% % - Placeholder for formal LaTeX pseudocode explaining LLaMEA's (1+1) Evolution Strategy parent-child swapping criteria.
% % - Placeholder to explicitly document the "Deterministic Ceiling" gap, justifying why clean-trained models fail when dropped into non-zero noise levels.
%
% ==============================================================================